\begin{lemma}[Bridge response transition]
Assume $F$ is a front on the full base and
$\mathsf{FrontBridge}(F,s,t,u)$.  Choose the next left node by
\[
 (s\in F\text{ and }s'=s)\quad\text{or}\quad
 (s\notin F\text{ and }s'=u).
\]
Suppose the right response is one of the following:
\begin{enumerate}
\item $t\notin F$, $t'$ is an ordered proper one-point extension of $t$
      in $T(F)$, and for its new point $n$ we set
      $u'=u\cup\{n\}$; or
\item $t\in F$, $t'=t$, and we choose a fresh $k$ larger than every member
      of $u$, setting $u'=u\cup\{k\}$.
\end{enumerate}
Then $\mathsf{FrontBridge}(F,s',t',u')$ holds.  The selected left node is
legal, and it is equal to $s$ when $s$ is terminal and a proper one-point
extension of $s$ otherwise.
\end{lemma}
\begin{proof}
The left legality is \noderef{FrontBridgeLeftMove}.  We check the two
right cases.  In the first case write $t'=t\cup\{n\}$.  Since $t$ is the
tail of $u$ whenever it is nonterminal, and the extension is ordered, the
tail of $u'=u\cup\{n\}$ is exactly $t'$.  If $t'$ is nonterminal this is
the exact-tail clause; if it is terminal it implies the required
initial-segment clause.  The old left relation to $u$ persists after the
new largest point is added.  If the selected left node is nonterminal, it
is $u$, and the ordered-extension conclusion
\noderef{FrontTreeOrderedExtension} puts $u'$ in $T(F)$; if it is
terminal, its old initial-segment relation to $u$ persists to $u'$.

In the second case $t'=t$ is terminal.  The new point $k$ is larger than
every member of $u$, so the tail of $u'$ is the old tail of $u$ followed by
$k$.  Thus the old terminal initial-segment relation for $t$ persists.
The same left-side argument applies: a nonterminal selected left node is
the old $u$ and has the ordered tree extension $u'$, while a terminal one
is an initial segment of $u'$.

The tree-membership clauses for $s'$ and $t'$ are supplied respectively by
\noderef{FrontBridgeLeftMove} and the response hypothesis.  Combining these
facts with the nonempty bridge $u'$ and the four defining clauses of
\noderef{FrontBridge} proves the transition.
\end{proof}
